% Generated draft: requires mathematical and source review.
\documentclass[11pt]{article}
\usepackage[margin=25mm]{geometry}
\usepackage{fontspec}
\setmainfont{Latin Modern Roman}
\usepackage{amsmath,amssymb,mathtools,unicode-math}
\setmathfont{Latin Modern Math}
\usepackage{enumitem,needspace,xurl,hyperref}
\setlength{\parindent}{0pt}
\setlength{\parskip}{5pt}
\setlength{\emergencystretch}{3em}
\newcommand{\N}{\mathbb{N}}
\newcommand{\Z}{\mathbb{Z}}
\newcommand{\Q}{\mathbb{Q}}
\newcommand{\R}{\mathbb{R}}
\newcommand{\C}{\mathbb{C}}
\newcommand{\F}{\mathbb{F}}
\newcommand{\A}{\mathbb{A}}
\newcommand{\PP}{\mathbb{P}}
\newcommand{\E}{\mathbb{E}}
\newcommand{\eps}{\varepsilon}
\newcommand{\dd}{\,\mathrm{d}}
\begin{document}
\section*{UnsolvedMath ID: 30000616}
\textbf{Source problem code:} catalogue entry
\subsection{Definitions and mathematical statement}
Boffa's conjecture, as attributed in the supplied original wording, asks for a model \(\mathcal M=(M,E)\models\mathrm{ZF}\), an element \(\kappa\in M\) with \(\mathcal M\models\text{``}\kappa\text{ is a cardinal''}\), and a nonidentity automorphism \(\sigma:M\to M\) satisfying
\[
xEy\quad\Longleftrightarrow\quad\sigma(x)E\sigma(y)
\qquad(x,y\in M),
\]
and
\[
\mathcal M\models 2^\kappa=\sigma(\kappa).
\]
The automorphism is a bijection of the model's domain preserving membership. It is not required to be an element or an internally definable class of the model.

The exponentiation and cardinal equality are evaluated \emph{inside} \(\mathcal M\), not in the surrounding universe. Since choice is not assumed, the ordinal-cardinal reading can be made precise by requiring an internally existing bijection between the internal power set \(\mathcal P^{\mathcal M}(\kappa)\) and the ordinal cardinal \(\sigma(\kappa)\). This also requires that particular internal power set to be well-orderable; it does not assume global choice. A different representation of cardinals in ZF would need to be fixed explicitly.

No external well-foundedness or transitivity of \(\mathcal M\) is required in the source. A well-founded extensional membership structure has no nonidentity membership automorphism, so imposing such a condition would change the question fundamentally. Internal satisfaction of Foundation is compatible with using a model that is ill-founded externally. The equality concerns one specified cardinal, not an assertion that \(\sigma\) acts by exponentiation on every cardinal.

The source does not give a countability or set-versus-class-model convention or an explicit relative-consistency assumption. In the usual model-existence interpretation these metamathematical choices should be stated before a consistency theorem is claimed. The supplied literature triage checked 2026-08-21 reports no authoritative complete resolution of the extracted claim; that is a dated assessment, not an independent consistency result.
\subsection{Short English statement}
Seek a ZF model with a membership automorphism sending one internal cardinal to the internal cardinality of its power set. All cardinal arithmetic is internal, and the model must not be assumed externally well founded or transitive.
\subsection{Sources}
\begin{itemize}
\item Local UnsolvedMath registry, permanent ID 30000616, catalogue entry, \emph{Boffa's Set-Theoretic Automorphism Conjecture}; the original wording attributes it to Boffa.
\item \emph{Mini-workshop: Logic, Combinatorics and Independence Results} (2007), Oberwolfach Reports, pp. 3087--3130, DOI \url{https://doi.org/10.4171/owr/2006/52}.
\item Supplied source-verification claim and literature triage checked 2026-08-21 cite only the report; no separate consistency theorem or cardinal-convention reference is supplied.
\end{itemize}
\end{document}
